\section{Methodology}
\label{sec:methodology}

Building on the observation that truthfulness and robustness may share underlying mechanisms, we design a cross-benchmark correlation study with explicit mechanism testing. Our approach addresses three questions: (1) Does a significant correlation exist after controlling for model size? (2) Does calibration correlate with either metric? (3) Does calibration moderate the correlation strength?

\subsection{Overview}

We evaluate decoder-only LLMs on three benchmarks: TruthfulQA MC1 (truthfulness), AdvGLUE (robustness), and MMLU (calibration reference). We compute partial correlations controlling for model size (log-transformed parameter count) and test whether Expected Calibration Error (ECE) explains the relationship.

\subsection{Benchmark Selection}

\textbf{TruthfulQA MC1} \cite{lin2022truthfulqa}: Multiple-choice format measuring whether models select truthful answers over common misconceptions. We use MC1 (single correct answer) for cleaner accuracy measurement.

\emph{Rationale:} TruthfulQA specifically measures factual accuracy under adversarial questioning designed to elicit falsehoods, making it complementary to AdvGLUE's perturbation-based robustness.

\textbf{AdvGLUE} \cite{wang2022advglue}: Adversarial variants of GLUE tasks with word-level perturbations. We report average accuracy across subtasks.

\emph{Rationale:} AdvGLUE tests robustness to input perturbations that preserve semantics but fool models, representing a different trust dimension from truthfulness.

\textbf{MMLU} \cite{hendrycks2021mmlu}: Multi-task academic benchmark. Used only for ECE computation on a neutral task.

\emph{Rationale:} Computing ECE on TruthfulQA or AdvGLUE would conflate calibration measurement with the metrics being correlated. MMLU provides a neutral reference.

\subsection{Model Selection}

We evaluate models from four families to avoid family-specific artifacts:

\begin{table}[h]
\centering
\begin{tabular}{lll}
\toprule
\textbf{Family} & \textbf{Models} & \textbf{Size Range} \\
\midrule
Pythia & 70M, 160M, 410M, 1B, 1.4B, 2.8B, 6.9B, 12B & 70M--12B \\
Llama-2 & 7B, 13B, 70B (base) & 7B--70B \\
Mistral & 7B & 7B \\
Falcon & 7B, 40B & 7B--40B \\
\bottomrule
\end{tabular}
\end{table}

\emph{Rationale:} Multiple families across a range of scales (70M--70B) ensure correlations are not artifacts of specific architectures or training procedures.

\subsection{Statistical Methods}

\subsubsection{Partial Correlation}

We compute Pearson partial correlation between TruthfulQA MC1 and AdvGLUE accuracy, controlling for $\log(\text{parameters})$:

\begin{equation}
r_{xy \cdot z} = \frac{r_{xy} - r_{xz} \cdot r_{yz}}{\sqrt{(1 - r_{xz}^2)(1 - r_{yz}^2)}}
\end{equation}

\emph{Rationale:} Larger models tend to perform better on both metrics. Without controlling for size, we might observe a spurious correlation driven by scale alone.

\subsubsection{Bootstrap Confidence Intervals}

We compute 95\% confidence intervals via 1000 bootstrap iterations, resampling models with replacement and recomputing partial correlation.

\emph{Rationale:} With $N=14$ models, asymptotic confidence intervals may be unreliable. Bootstrap provides nonparametric uncertainty quantification.

\subsubsection{Expected Calibration Error (ECE)}

For each model, we compute 15-bin ECE on MMLU:

\begin{equation}
\text{ECE} = \sum_{b=1}^{B} \frac{|B_b|}{N} |\text{acc}(B_b) - \text{conf}(B_b)|
\end{equation}

where bins are formed by predicted probability and we measure the gap between accuracy and confidence.

\emph{Rationale:} ECE is the standard calibration metric. 15 bins balance granularity with stability given sample sizes.

\subsubsection{Mechanism Testing}

\textbf{H-M1 (Calibration-Metric Correlation):} If calibration underlies both capabilities, ECE should negatively correlate with TruthfulQA and AdvGLUE (lower ECE = better calibration = higher performance).

\textbf{H-M2 (Calibration Moderation):} If calibration mediates the truthfulness-robustness relationship, well-calibrated (low-ECE) models should show stronger correlation than poorly-calibrated models. We test via Fisher's z-test comparing correlation coefficients across ECE tertiles.

\subsection{Hypothesis Structure}

Our analysis follows a hierarchical verification plan:

\begin{table}[h]
\centering
\begin{tabular}{llll}
\toprule
\textbf{ID} & \textbf{Type} & \textbf{Test} & \textbf{Gate} \\
\midrule
h-e1 & Existence & Partial $r > 0.3$, $p < 0.05$ & MUST\_WORK \\
h-m1 & Mechanism & ECE-metric $r < -0.2$ & SHOULD\_WORK \\
h-m2 & Mechanism & Low-ECE $r >$ High-ECE $r$ & SHOULD\_WORK \\
h-c1 & Condition & Pattern holds for base and IT separately & SHOULD\_WORK \\
\bottomrule
\end{tabular}
\end{table}

\emph{Rationale:} The MUST\_WORK gate (h-e1) determines whether the core correlation exists. SHOULD\_WORK gates test the calibration mechanism; failure redirects to alternative explanations rather than rejecting the main finding.

\subsection{Evaluation Protocol}

All evaluations use lm-evaluation-harness \cite{gao2023lmeval} with consistent settings:
\begin{itemize}
    \item Batch size: auto (determined by model size)
    \item Prompt format: default harness templates
    \item Sampling: greedy decoding for reproducibility
\end{itemize}

We prioritize reproducibility and fair comparison over optimizing individual model performance.
